\subsection{The same cover obstruction with a multiple inverse image}
\paragraph{A2Zr (Extension of the construction exclusion; submitted for audit).}
Let $\pi:X\to B$ be a finite flat generically separable morphism of
smooth projective integral varieties of dimension $n\geq2$ over an
algebraically closed field of characteristic $p>0$. Let $C\subset B$
be a smooth ample integral divisor. Suppose
\[
 \pi^*C=eY
\]
as Cartier divisors, where $e\geq1$ and $Y$ is smooth integral, and
suppose $\rho=\pi|_Y:Y\to C$ has degree $p$.
Then no nontrivial $M\in\operatorname{Pic}(C)[p](k)$ is trivialized
by $\rho$. In particular, increasing the ramification index in the
A2Z construction cannot evade that exclusion.

Assume such $M$ exists. Exactly as in A2Z, the compatible
trivialization defines a finite birational map from $Y$ to the
integral $\mu_p$-torsor of $M$. Thus $Y$ is its normalization,
$\rho$ is purely inseparable, and finite duality gives
\[
 \omega_Y\simeq\rho^*\omega_C\otimes\mathcal O_Y(-D)
                                                               \tag{A2Zr.1}
\]
for an effective conductor divisor $D$. No normality assumption on
the torsor is made.

Let $R$ be the effective divisor of the determinant of
$d\pi:\pi^*\Omega_B^1\to\Omega_X^1$. We claim that
\[
 \operatorname{ord}_Y(R)\geq e.                                \tag{A2Zr.2}
\]
This can be checked at the generic point of $Y$. After shrinking
around the generic points, take a parameter $t$ for $C$ and functions
$x_2,\ldots,x_n$ whose differentials restrict to a basis of
$\Omega^1_{k(C)/k}$. Then $dt,dx_2,\ldots,dx_n$ form a basis of
$\Omega^1_B$ there. Similarly take a parameter $u$ for $Y$ and a
basis of residue differentials on $Y$, so $du$ is part of a basis
of $\Omega^1_X$. The equality of Cartier divisors gives
$\pi^*t=v u^e$ for a unit $v$.

If $p\mid e$, then
\[
 d(\pi^*t)=u^e\,dv,
\]
so the determinant is divisible by $u^e$. If $p\nmid e$, its possible
term of order $e-1$ is a unit times
\[
 u^{e-1}\,du\wedge d(\pi^*x_2)\wedge\cdots\wedge d(\pi^*x_n).
\]
Modulo $u$, the last $n-1$ factors projected to
$\Omega^1_{k(Y)/k}$ are the differential of the inseparable map
$\rho$. Their top exterior product is zero, since a finite
inseparable map between smooth varieties of equal dimension has
zero determinant on differentials. Therefore this term too is
divisible by $u^e$, proving (A2Zr.2).

Write $R=eY+R'$ with $R'\geq0$. The canonical bundle formula and
adjunction, keeping the multiplicity $e$, give
\begin{align*}
 \omega_Y\otimes\rho^*\omega_C^{-1}
 &\simeq\mathcal O_X(R+Y-\pi^*C)|_Y\\
 &\simeq\mathcal O_X(R-(e-1)Y)|_Y
  =\mathcal O_X(Y+R')|_Y.
\end{align*}
The divisor $Y$ is ample because $eY=\pi^*C$ is ample. Comparing
with (A2Zr.1) and taking degree against
$\mathcal O_Y(Y)^{n-2}$ yields
\[
 0=Y^n+R'\cdot Y^{n-1}+D\cdot(Y|_Y)^{n-2}>0,
\]
a contradiction. This proves the extension.

The residue degree $p$ and the absence of other components in
$\pi^*C$ are explicit hypotheses. If additional components occur,
adjunction subtracts their restrictions to $Y$; if the residue
degree is larger, $Y$ need not be the normalization of the single
$\mu_p$-torsor. Neither situation is decided by this argument.
